% engines: lualatex
% expect-warnings: 0
% expect-errors: 0
\documentclass{article}
\usepackage[margin=1cm]{geometry}
\usepackage[kernel = lua]{../../latex/pgf-go}
\newcommand\case[1]{\par\medskip\noindent\textbf{#1}\par\nopagebreak}
\pagestyle{empty}

% Draws frames of \sgfanimation one by one (as in sgf-animation.tex).
\ExplSyntaxOn
\NewDocumentCommand \frames {mmm} {
	\group_begin:
		\pgfgo_sgfanimation_setup:nn {#1} {#2}
		\clist_map_inline:nn {#3} {\pgfgo_sgfanimation_frame:n {##1}\quad}
	\group_end:
}
\ExplSyntaxOff
\begin{document}

% joseki.sgf: 19x19, upper right corner. B R16, W P17; then B N17, W R14, B R12 (main line) or
% B Q15, W N17 with the label a at O16. The whole tree uses N12:R17.

\case{goban with captures (G9 is captured) and with \string\sgf}
\begin{goban}[19, partial = auto, captures]
	\stones[b]{F9, G10, H9}\stones[w]{G9}\moves{G8}
\end{goban}
\begin{goban}[19, partial = auto]\sgf[to = 2]{joseki.sgf}\end{goban}
\begin{goban}[19, partial = auto]\sgf{joseki.sgf}\end{goban}

\case{sgf diagrams, partial = auto: the same part (M11:T19) in every diagram, also in the variation}
\sgfdiagrams[partial = auto, step = 2, variations, goban = {scale = 0.5}]{joseki.sgf}

\case{sgf diagrams, partial = auto with goban = \{partial margin = 0\}: N12:T19}
\sgfdiagrams[partial = auto, goban = {scale = 0.5, partial margin = 0}]{joseki.sgf}

\case{sgf diagrams: a view wins over partial = auto (A5:E9, then the whole board, then F1:I4)}
\sgfdiagrams[partial = auto, step = 2, goban = {scale = 0.4}]{views.sgf}

\case{sgf animation, partial = auto: frames 0, 2 and 5 have the same board}
\frames{partial = auto, goban = {scale = 0.4}}{joseki.sgf}{0, 2, 5}

\end{document}
